\documentclass[a4paper]{article}
\usepackage{amsfonts}
\usepackage{amsmath}
\usepackage{amssymb}
\usepackage{graphicx}     

\newcommand{\dint}{\displaystyle\int}

\begin{document}
  
\noindent \large Auteur: Thijs Van Schel \\
\noindent \large Studierichting: Informatica\\
\noindent \large Oefeningenreeks 4A nr. 42.18\\

\medskip

\normalsize

$\dint \dfrac{2x^2-5x+7}{x^3-7x+6}dx$\\

$\begin{array}{c|ccc|c}
  & 1 & 0 & -7 & 6 \\
1 & & 1 & 1 & -6 \\
\hline
  & 1 & 1 & -6 & 0
\end{array}$\\

$\Rightarrow (x-1)(x^2+x-6) = (x-1)(x-2)(x+3)$\\

$\dint \dfrac{2x^2-5x+7}{(x-1)(x-2)(x+3)}dx$\\

$\dfrac{2x^2-5x+7}{(x-1)(x-2)(x+3)}=\dfrac{A}{x-1}+\dfrac{B}{x-2}+\dfrac{C}{x+3}$\\

$A = \left . \dfrac{2x^2-5x+7}{(x-2)(x+3)}\right |_{x=1}=-1$ \\

$B = \left . \dfrac{2x^2-5x+7}{(x-1)(x+3)}\right |_{x=2}=1$\\

$C = \left . \dfrac{2x^2-5x+7}{(x-1)(x-2)} \right |_{x=-3}=2$\\

$\Rightarrow -\dint \dfrac{1}{(x-1)}dx + \dint \dfrac{1}{(x-2)}dx + 2\dint \dfrac{1}{(x+3)}dx$\\

$= -\ln|x-1|+\ln|x-2|+2\ln|x+3|+C$

\begin{thebibliography}{999}
%\bibitem{a} Referentie
\end{thebibliography}
\end{document} 